\chapter{Endoscopic Retrograde Cholangiopancreatography}

\section*{Overview}
Endoscopic retrograde cholangiopancreatography (ERCP) is an endoscopic procedure that examines and treats the bile ducts, gallbladder, and pancreatic ducts.

\section*{Alternative Names}
ERCP; endoscopic cholangiopancreatography; biliary endoscopy

\section*{Principle}
An endoscope is passed to the duodenum, and the ampulla of Vater is cannulated. Contrast is injected to image the biliary and pancreatic ducts, and instruments can open, dilate, or stent these ducts and retrieve stones.

\section*{History}
ERCP was developed in the late 1960s and 1970s, initially as a diagnostic tool, and evolved into a primarily therapeutic procedure after endoscopic sphincterotomy was introduced in 1974.

\section*{Why It Is Done}
Performed for obstructive jaundice, choledocholithiasis, cholangitis, bile leaks, and pancreatic disease such as strictures, pseudocysts, and some tumors, and for palliation of malignant obstruction with stents.

\section*{Is This a Primary or Secondary Test or Procedure}
Primary therapeutic procedure for biliary and pancreatic obstruction.

\section*{How It Is Performed}
Performed under conscious sedation or general anesthesia. The endoscope reaches the duodenum, the ampulla is cannulated, contrast outlines the ducts under fluoroscopy, and interventions such as sphincterotomy, stone extraction, dilation, or stenting are performed.

\section*{Preparation}
NPO before the procedure, IV fluids, and pre-procedure labs including liver tests and coagulation studies. Antibiotics are given in selected cases.

\section*{Contraindications and Precautions}
Severe acute pancreatitis without obstruction, coagulopathy, and inability to tolerate sedation are relative contraindications. Precaution: post-ERCP pancreatitis is the most common complication.

\section*{Risks}
Post-ERCP pancreatitis, bleeding, perforation, cholangitis, sedation-related complications, and adverse reactions to contrast.

\section*{Aftercare and Recovery}
The patient is observed for signs of pancreatitis or bleeding; blood tests and imaging may be checked the next day.

\section*{Outcome and Follow-up}
Findings may show stones, strictures, leaks, or malignancy. Successful sphincterotomy and drainage are confirmed by contrast and symptom improvement.

\section*{Expected Results}
A patent, non-obstructed biliary and pancreatic ductal system with no stones, stricture, or leak.

\section*{Follow-up}
Liver function tests and imaging follow-up; repeat ERCP may be needed for stent exchange or recurrent stones.

\section*{References}
\begin{enumerate}
  \item MedlinePlus. Endoscopic retrograde cholangiopancreatography (ERCP). U.S. National Library of Medicine; 2025.
  \item Current Medical Diagnosis and Treatment. McGraw-Hill; 2025.
\end{enumerate}

\clearpage
